% =====================================================================
%  bg_A_knowledge_graphs.tex
%  WHERE:  inside Section 2.2 "Knowledge Graphs",
%          directly AFTER your "Example 1" box (KG_1 = (N_1, E_1, L_1))
%          and BEFORE  \section{Language Models}.
%  ADDS:   2.2.1 SPARQL, 2.2.2 Open-World Assumption, 2.2.3 KG Enrichment
% =====================================================================

\subsection{Querying Knowledge Graphs with SPARQL}
\label{sec:bg-sparql}

Once knowledge is stored as RDF triples, a dedicated language is needed to
retrieve it. The \emph{SPARQL Protocol and RDF Query Language} (SPARQL) is
the W3C standard for querying RDF data~\cite{harris2013sparql}. A SPARQL query
describes a small graph pattern, that is, a set of triples in which some
positions are replaced by variables (written with a leading \texttt{?}). The
query engine then searches the Knowledge Graph for all sub-graphs that match
this pattern and returns the values bound to the variables. In addition to
simple look-ups, SPARQL supports filtering (\texttt{FILTER}), optional
patterns (\texttt{OPTIONAL}), aggregation (\texttt{COUNT}, \texttt{GROUP BY})
and federated queries over several endpoints (\texttt{SERVICE}). SPARQL is
therefore the natural interface between a Knowledge Graph and the
applications built on top of it, including the data preparation steps for
language models described later in this chapter.

\begin{bgdefinition}{Basic Graph Pattern}
Let $V$ be an infinite set of variables disjoint from $Con$. A \emph{triple
pattern} is a tuple $t \in (Con \cup V) \times (Con \cup V) \times (Con \cup V)$,
and a \emph{basic graph pattern} $P$ is a finite set of triple patterns.
A \emph{solution} of $P$ over a Knowledge Graph $KG=(N,E,L)$ is a mapping
$\mu : V \rightarrow Con$ such that replacing every variable $v$ in $P$ by
$\mu(v)$ yields a set of triples that is contained in $E$, i.e.\
$\mu(P) \subseteq E$~\cite{harris2013sparql,hogan2021knowledge}.
\end{bgdefinition}

\begin{bgexample}[SPARQL over the questionnaire Knowledge Graph]
In the DigiStrucMed questionnaire Knowledge Graph, every patient answer is
linked to the question it answers and to an answer category (e.g.\
\emph{very\_positive}, \emph{neutral}, \emph{negative}). The following query
counts, for every question of the category \emph{Quality}, how many
participants selected each answer category.
\begin{lstlisting}[style=bgcode]
PREFIX ds: <http://digistrucmed.org/questionnaire#>
SELECT ?questionDesc ?category (COUNT(DISTINCT ?resp) AS ?numPeople)
WHERE {
  ?question a ds:Question ;
            ds:hasQuestionCategory ds:Quality ;
            ds:QuestionDescription ?questionDesc .
  ?resp     a ds:Answered ;
            ds:hasAnswer ?answer .
  ?answer   ds:isAnswerTo ?question ;
            ds:hasAnswerCategory ?category .
}
GROUP BY ?questionDesc ?category
ORDER BY ?questionDesc ?category
\end{lstlisting}
The \texttt{WHERE} clause is a basic graph pattern with seven triple
patterns; each result row is one solution $\mu$, aggregated by
\texttt{GROUP BY}.
\end{bgexample}

\subsection{Open-World Assumption and Knowledge Graph Incompleteness}
\label{sec:bg-owa}

Knowledge Graphs are almost never complete: guidelines are updated, sources
disagree, and automatic extraction misses facts. How a system interprets a
\emph{missing} triple is therefore an important design decision. Relational
databases usually follow the \emph{Closed-World Assumption} (CWA), under
which every fact that is not stored is considered false. Semantic Web
standards such as RDF and OWL instead follow the \emph{Open-World Assumption}
(OWA)~\cite{hogan2021knowledge}.

\begin{bgdefinition}{Open-World Assumption}
Under the Open-World Assumption, a triple $\langle s, p, o\rangle \notin E$
is \emph{unknown} rather than false: the Knowledge Graph
$KG=(N,E,L)$ is regarded as an incomplete description of the domain, and the
absence of a triple neither confirms nor refutes the corresponding
statement~\cite{hogan2021knowledge}.
\end{bgdefinition}

\begin{bgexample}[Missing is not false]
Assume a guideline Knowledge Graph contains the triple
\texttt{(sacubitril/valsartan, recommendedFor, HFrEF)}, but no triple that
connects sacubitril/valsartan to arterial hypertension. Under the CWA, a system would conclude that
the drug is \emph{not} related to hypertension. Under the OWA, the system can
only conclude that the graph \emph{does not state} such a relation. For a
patient-facing assistant this difference matters: an answer such as
``this drug has nothing to do with blood pressure'' would be an unjustified
claim derived from incomplete data.
\end{bgexample}

The OWA has two direct consequences for this thesis. First, it motivates
\emph{reasoning} over Knowledge Graphs, i.e.\ deriving facts that are not
stated explicitly but follow from other facts (Section~\ref{sec:bg-nesy}).
Second, it complicates the evaluation of such reasoning, because a predicted
fact that is missing from the graph is not necessarily wrong. Rule-mining
systems therefore introduce intermediate assumptions, such as the
Partial Completeness Assumption discussed in
Section~\ref{sec:bg-rule-mining}.

\subsection{Knowledge Graph Enrichment with External Resources}
\label{sec:bg-enrichment}

A Knowledge Graph that is built from a single type of source, for example the
text of clinical guidelines, typically contains surface forms such as
abbreviations (\emph{MRA}, \emph{SGLT2i}) or spelling variants
(\emph{beta-blocker}, \emph{beta blockers}) without any further description.
\emph{Knowledge Graph enrichment} links such entities to external,
curated resources and imports additional facts about them. In the
biomedical domain, three resources are particularly relevant:

\begin{itemize}
  \item \textbf{Wikidata}, a free, collaboratively edited, multilingual
        Knowledge Graph that provides structured statements, synonyms
        (aliases) and identifiers for millions of entities, including
        diseases and drugs, through a public SPARQL endpoint~\cite{vrandecic2014wikidata}.
  \item \textbf{The Unified Medical Language System (UMLS)}, which integrates
        more than one hundred biomedical vocabularies (e.g.\ SNOMED~CT,
        MeSH, ICD) under shared concept identifiers~\cite{bodenreider2004umls}.
  \item \textbf{DrugBank}, a curated knowledge base that combines detailed
        chemical, pharmacological and drug-target information~\cite{wishart2018drugbank}.
\end{itemize}

\begin{bgexample}[Enriching a guideline abbreviation]
The guideline Knowledge Graph of this thesis stores the abbreviation
\emph{MRA} as a surface form extracted from guideline text. Enrichment adds
its expansion and links the drug class to a Wikidata entity:
\begin{lstlisting}[style=bgcode]
ds:entity/MRA  a ds:Abbreviation, ds:SurfaceForm ;
    ds:hasExpansion ds:entity/Mineralocorticoid%20receptor%20antagonist .

ds:entity/beta%20blocker  a ds:WikidataEntity, ds:SurfaceForm, ds:Drug ;
    ds:subclass_of ds:entity/medication .
\end{lstlisting}
Because all pipelines mint IRIs in the same namespace, a guideline mention
and a Wikidata concept with identical text become \emph{one} node. This is
what allows facts from different sources to be combined in a single query.
\end{bgexample}

Enrichment increases the coverage and connectivity of the graph, which is
important both for answering patient questions and for mining rules that
generalise across documents. At the same time, imported facts inherit the
quality and the licence of their source, and automatic entity linking can
introduce errors. Enriched facts should therefore keep their provenance
(e.g.\ the source document or the Wikidata identifier), so that they can be
traced back and verified by medical experts.
